% !TeX root = ../papers/nc-p3.tex
\section{台144の楕円K3曲面と水平三層の無印商}
\label{sec:parallel-horizontal-moduli}
\begin{theorem}\label{thm:parallel-horizontal-moduli}
112型の主成分で、台144,146,147の幾何学的一般点は標準Koszul左右AS正則で、最小周期は4である。
指定始点を固定した無印同型類の関数体は順に次の通りである。
\begin{enumerate}
\item 台144ではPicard数20、Néron--Severi判別式の絶対値8の楕円K3曲面の関数体である。
\item 台146では $k(t,b,a)$、台147では四変数の純超越拡大である。
\end{enumerate}
ここで曲面・三次元多様体は指定した台の層の双有理模型を表す。
\end{theorem}

\begin{proof}[曲面と三次元層]
前節までの $B=1$, $a=A/t$, $y=C$, $m=t^2=b+c-bc$ を使う。
台144では $y=1-a$ かつ $c=b^2/E$, $E=(b-a)(b-1+a)$ である。
$t^2=b a(1-a)/E$ の右辺は $b=0$ で奇数位数を持つから二重被覆は既約である。
$W=tE$ と置くと $W^2=b a(1-a)E$ であり、$b$ の関数体上で
\[
 k_0=-b^2(1-b),\qquad X=k_0/a-k_0,\qquad Z=k_0 W/a^2
\]
により $Z^2=X(X+L)(X+M)$ となる。ただし $L=b(1-b)^2$, $M=b^3$ である。
逆写像は $a=k_0/(X+k_0)$ から得られる。判別式 $16L^2M^2(L-M)^2$ と
$c_4=16(L^2-LM+M^2)$ の位数を調べると
\[
\begin{array}{c|rrrr}
 b&0&1&1/2&\infty\\\hline
 (v(c_4),v(\Delta))&(2,10)&(0,4)&(0,2)&(2,8)\\
 \text{Kodaira型}&I_4^*&I_4&I_2&I_2^*
\end{array}
\]
となる。無限遠では $X/b^4,Z/b^6$ の整モデルを使い、全ての点で $v(\Delta)<12$ なので極小である。
他のファイバーは滑らかで、基本直線束は $\mathcal O_{\mathbf P^1}(2)$。
切断付き楕円面の標準束公式と $R^1f_*\mathcal O=\mathcal O(-2)$ から極小滑らか模型はK3である。
可約ファイバーの根格子 $D_8,A_3,A_1,D_6$ から自明格子の階数は20となり、
標数零のK3のPicard数の上界およびShioda--Tate公式からMordell--Weil階数は0である。
加法的ファイバーの連結な滑らかな部分と還元核には標数零で有限位数点がないため、
捩れ部分群は $I_4^*$ の位数4の成分群に単射する。三つの非自明な有理2捩れ点を既に持つので位数は4であり、
判別式の絶対値は $4\cdot4\cdot2\cdot4/4^2=8$ となる。

台146では $t^2(a-b)(y-b)=b(a-1)(y-1)$ である。$y$ についての係数と定数項は
$k[t,b,a]$ で互いに素であり、全体が既約で、
\[
 y=\frac{b\big(t^2(a-b)-(a-1)\big)}{t^2(a-b)-b(a-1)} .
\]
従って $b=1,t^2=1$ における余分な自由度から別の成分を生じることもない。
元の各台では薄い係数・行列式・指定した厚い係数の非零条件を課す。
\end{proof}

\begin{proof}[文字枠と無印同型]
三層の閉包はいずれも台143上の $t=1,y=2,a=-1,b=-2,c=1$ を含む。
この点は正則で、同じ二切断の階数を持つ。定理\ref{thm:parallel-invertible-rectangles}の次数別の特殊化論証で
有限Frobenius双対および各内部次数でのKoszul性を幾何学的一般点に移せる。
この点では位置1,3の二つの部分traceおよびその固有線を挿入した位置0,2の縮約が四枠を復元する。
各々に要る固有値差の非消滅は開条件なので、三層の一般点でも枠は内在的である。

台144,146で指定始点の台を保つ四置換は恒等のみで、薄い成分の規格化と $B=1$ により全係数が一意である。
奇数移動の台置換はなく、二位置移動の台置換は全位置で $(0\,1)$ の一つだけである。
係数には $(a,y)\mapsto(y,a)$ と作用する。このため台144で二位置移動が $a\mapsto1-a$ になる一方、
$t\mapsto-t$ は別の一般同型類を与える。台146でも二位置移動は一般点を固定しない。

台147は $b\ne1$ で $c=(m-b)/(1-b)$ と書ける。指定始点の全作用は恒等と
\[
 I(t,b,c,a,y)=(t,m/b,m/c,1-c+cy/b,1-c+ca/b)
\]
である。一位置移動の代表は
\[
 J(t,b,c,a,y)=(1/t,c/m,1/b,c(1-a)/m,c(1-y)/m).
\]
$J$ は $I$ と可換で、$J^2$ は $(a,y)$ の交換と $I$ の合成である。
これらを確かめるには台を移す置換を全探索し、薄い係数の $t$ の指数について位置別倍率を消去する。
残る倍率で $B=1$ に戻した後、16個の厚い係数の全てが表示先の補完式に一致する。
一般には $t\ne1/t$ かつ $(a,y)$ の交換が $\{1,I\}$ の同じ軌道に入らず、周期は4である。

$u=(b-t)/(b+t)$ と置くと $I(u)=-u$ なので、$A_0=(a+I(a))/2$,
$A_1=(a-I(a))/(2u)$ は不変である。$b=t(1+u)/(1-u)$, $a=A_0+u A_1$ および
$y=(A_0-u A_1-1+c)b/c$ で全て復元される。従って不変体は $k(t,u^2,A_0,A_1)$ である。
\path{computations/verify_parallel_horizontal_moduli.py} はファイバーの位数、消去式、台の全置換と全係数の変換を検算する。
\end{proof}
